\chapter{Lice}
\index{Lice}

\section*{Definition and Overview}
Pediculosis is an infestation of the skin or hair by tiny, wingless, parasitic insects known as lice (singular: louse) that feed on human blood. The three primary species that infest humans are \textit{Pediculus humanus capitis} (head lice), \textit{Pediculus humanus corporis} (body lice), and \textit{Pthirus pubis} (pubic lice, or crabs). While head and pubic lice cause localised pruritus and discomfort, body lice are unique as vectors for severe systemic pathogens. The infestation is highly contagious and spread through close contact or shared personal items, representing a significant public health issue globally.

\section*{Epidemiology}
Lice infestations occur worldwide, affecting individuals of all socioeconomic backgrounds. Head lice are particularly prevalent among children aged 3 to 11 years in child care and school settings, with millions of cases reported annually. Girls are affected more frequently than boys, likely due to hair length and social play patterns. In contrast, body lice infestations are typically restricted to populations experiencing severe poverty, homelessness, overcrowding, or poor hygiene conditions where regular laundering of clothing is unavailable. Pubic lice are most common in adolescents and young adults, primarily transmitted via sexual contact, with a global prevalence estimated at 1\% to 2\% of the population.

\section*{Aetiology and Pathophysiology}
The pathology of pediculosis is driven by the life cycle of the lice and the host's immune response to the parasite's salivary secretions. Each species has specific habitat preferences on the human body.
\begin{itemize}
    \item \textbf{Head lice (\textit{Pediculus humanus capitis}):} These reside on the scalp, particularly behind the ears and near the nape of the neck. They cement their eggs (nits) to the hair shaft close to the scalp surface, where temperature and humidity support incubation.
    \item \textbf{Body lice (\textit{Pediculus humanus corporis}):} Unlike head lice, body lice reside in the seams of clothing and bedding, only migrating to the skin to feed. They are associated with poor hygiene and lack of laundering.
    \item \textbf{Pubic lice (\textit{Pthirus pubis}):} These infest the coarse hair of the pubic region, but can also spread to the eyelashes, eyebrows, beard, and axillary hair, especially in highly infested individuals.
    \item \textbf{Hypersensitivity reaction:} The primary symptom, pruritus, is not caused by the bite itself but by an inflammatory immunological reaction to the louse's saliva, which contains vasodilatory and anticoagulant substances injected during feeding.
    \item \textbf{Vector transmission:} Body lice are unique in their ability to transmit life-threatening bacterial infections, including epidemic typhus (\textit{Rickettsia prowazekii}), trench fever (\textit{Bartonella quintana}), and louse-borne relapsing fever (\textit{Borrelia recurrentis}). Head and pubic lice do not transmit these pathogens.
\end{itemize}

\section*{Clinical Features}
The clinical presentation of pediculosis varies depending on the type of louse, but pruritus is the hallmark feature across all forms.
\begin{itemize}
    \item \textbf{Intense pruritus:} Severe itching of the scalp, body, or pubic region, which is typically worse at night when the parasites are most active.
    \item \textbf{Visible lice and nits:} Detection of live insects (1 to 3~mm in length) crawling on the scalp, skin, or clothing, or small, oval, white-to-yellow eggs (nits) firmly glued to hair shafts or clothing fibres.
    \item \textbf{Excoriations and secondary infection:} Scratches and skin abrasions from constant scratching, which can become infected by bacteria (such as \textit{Staphylococcus aureus} or \textit{Streptococcus pyogenes}), leading to impetigo, oozing, and local lymphadenopathy.
    \item \textbf{Maculae caeruleae:} Characteristic slate-grey or bluish macules (typically 0.5 to 2~cm in size) on the skin of the trunk or thighs in pubic lice infestations, caused by the breakdown of red blood cells by the louse salivary enzymes.
    \item \textbf{Eyelash involvement (Phthiriasis palpebrarum):} Burning, itching, and crusting of the eyelids, with visible nits and lice on the eyelashes, primarily seen in children with pubic lice infestations.
    \item \textbf{Vagabond's disease:} Hyperpigmentation and thickening (lichenification) of the skin in individuals with long-standing, untreated body lice infestations.
\end{itemize}

\section*{Differential Diagnosis}
To ensure appropriate treatment, pediculosis must be distinguished from other dermatological conditions that cause itching and scale-like lesions.
\begin{itemize}
    \item \textbf{Seborrheic dermatitis:} Causes scaling and itching of the scalp, but the flakes (dandruff) are greasy, diffuse, and easily brushed out of the hair, unlike nits which are firmly attached.
    \item \textbf{Scabies:} An infestation by the \textit{Sarcoptes scabiei} mite that causes intense nocturnal itching and burrows, typically in the web spaces of the fingers, wrists, and waistline, rather than the scalp.
    \item \textbf{Psoriasis:} A chronic autoimmune skin condition causing thick, silvery-scaled plaques on the scalp and body, which may itch but does not involve parasites.
    \item \textbf{Hair casts (pseudonits):} Keratinous rings that encircle the hair shaft and can slide freely along it, unlike nits which are fixed in place.
    \item \textbf{Contact dermatitis:} Allergic reaction to hair products, soaps, or laundry detergents, presenting with pruritus and erythema but no visible parasites.
\end{itemize}

\section*{When to Seek Medical Care}
Consult a healthcare professional if you suspect a lice infestation and over-the-counter treatments fail, if there are signs of secondary bacterial infection, or if the infestation is present in infants.

\textbf{Contact your local emergency services for:}
\begin{itemize}
    \item Signs of severe systemic infection or sepsis arising from secondary bacterial skin infections, such as high fever, confusion, rapid breathing, or spreading red streaks from an excoriation.
    \item Severe allergic reactions (anaphylaxis) to chemical pediculicide treatments, presenting with difficulty breathing, facial swelling, or collapse.
\end{itemize}
Other reasons to consult a doctor include persistent scalp itching without visible lice, vaginal discharge or discomfort in adolescents with pubic lice, or eyelash involvement in children.

\section*{Diagnosis and Investigations}
Diagnosis is almost exclusively made by direct visualisation of a live louse or viable nits, with mechanical techniques significantly increasing diagnostic accuracy.
\begin{itemize}
    \item \textbf{Wet combing:} The gold standard method for detecting head lice, using a fine-toothed detection comb (teeth spacing 0.2 to 0.3~mm) on hair lubricated with conditioner. This is up to four times more effective than visual inspection alone.
    \item \textbf{Visual inspection:} Checking the hair and scalp under bright light, focusing on the nape of the neck and behind the ears, looking for moving lice or nits located within 6~mm of the scalp.
    \item \textbf{Microscopic examination:} Examining collected lice or hair shafts under a light microscope to confirm species identification and differentiate viable nits (containing embryos) from empty egg casings.
    \item \textbf{Clothing inspection:} Inspecting the seams of clothing and bedding for body lice and their eggs, particularly in individuals with poor hygiene or who are homeless.
    \item \textbf{STI screening:} Recommending a comprehensive sexually transmitted infection screen for any patient diagnosed with pubic lice, as co-transmission of other STIs is common.
\end{itemize}

\section*{Management}
Treatment of pediculosis involves eliminating the parasites from the host and, in the case of body lice, disinfecting the environment.
\begin{itemize}
    \item \textbf{Topical pediculicides:} Application of chemical agents such as permethrin (1\% or 5\%), pyrethrins, malathion, or benzyl alcohol to the affected areas. A second application is typically required 7 to 10 days later to kill newly hatched nymphs before they can reproduce.
    \item \textbf{Wet combing therapy:} A non-chemical approach involving systematic combing of wet, conditioned hair with a fine-toothed comb every 3 to 4 days for at least two weeks, suitable for infants, pregnant women, or cases with chemical resistance.
    \item \textbf{Environmental decontamination:} Washing clothing, bedding, and towels used by the infested person in hot water (at least 60 degrees Celsius) and drying on a high heat cycle, or sealing non-washable items in plastic bags for two weeks to starve the lice.
    \item \textbf{Body lice management:} Primarily achieved through improved personal hygiene, regular bathing, and washing of clothing and bedding, as chemical treatment of the skin is rarely necessary once clean clothes are provided.
    \item \textbf{Sexual partner notification:} Treating all sexual contacts of patients with pubic lice within the past month to prevent reinfection.
    \item \textbf{Oral ivermectin:} An oral antiparasitic agent reserved for refractory cases or mass outbreaks, administered in two doses spaced 7 to 10 days apart.
\end{itemize}

\section*{Complications}
\begin{itemize}
    \item \textbf{Secondary bacterial infection:} Impetigo, ecthyma, or cellulitis caused by \textit{Staphylococcus aureus} or \textit{Streptococcus pyogenes} entering through scratched skin.
    \item \textbf{Sleep disturbance and psychological distress:} Chronic pruritus leads to insomnia, which can cause fatigue and irritability, while the social stigma of infestation often leads to anxiety and embarrassment.
    \item \textbf{Lice-borne infectious diseases:} Transmission of epidemic typhus, trench fever, or relapsing fever by body lice in vulnerable populations.
    \item \textbf{Blepharitis and conjunctivitis:} Inflammatory complications of the eyes when pubic lice infest the eyelashes.
    \item \textbf{Contact dermatitis and chemical burns:} Skin irritation or hypersensitivity reactions from repeated or inappropriate use of chemical pediculicides.
\end{itemize}

\section*{Prognosis and Outcomes}
The prognosis for pediculosis is excellent with appropriate, completed treatment. Head and pubic lice infestations are highly curable, and symptoms typically resolve rapidly once the parasites are eradicated. Resistance to common chemical treatments is an increasing global challenge, but alternative formulations and physical removal methods remain highly effective. For body lice, prognosis is closely tied to access to clean clothing, laundry facilities, and improved sanitation.

\section*{Prevention and Self-Care}
\begin{itemize}
    \item \textbf{Avoid direct contact:} Do not engage in head-to-head or close body contact with individuals known to have a lice infestation.
    \item \textbf{Do not share personal items:} Avoid sharing combs, brushes, hats, scarves, towels, headphones, or bedding.
    \item \textbf{Regular screening:} Perform weekly wet combing checks on children during school outbreaks to detect and treat infestations early.
    \item \textbf{Proper hair hygiene:} Keep hair tied back or braided in school environments to reduce the risk of head lice transfer.
    \item \textbf{Laundering practices:} Regularly wash clothing and bedding at high temperatures, particularly when returning from group camps or travel.
    \item \textbf{Safe sexual practices:} Limit sexual contact with individuals who have active pubic lice infestations until they have been fully treated.
\end{itemize}

\section*{Sources and Further Reading}
The following reputable organisations provide further information and clinical guidelines on the prevention, diagnosis, and treatment of lice infestations:
\begin{itemize}
    \item Centers for Disease Control and Prevention (CDC). \textit{Parasites - Lice.} https://www.cdc.gov/
    \item Healthdirect Australia. \textit{Head lice.} https://www.healthdirect.gov.au/
    \item Better Health Channel. \textit{Head lice (pediculosis).} https://www.betterhealth.vic.gov.au/
    \item World Health Organization (WHO). \textit{Louse infestations.} https://www.who.int/
    \item National Health and Medical Research Council (NHMRC). \textit{Clinical guidelines for head lice control.} https://www.nhmrc.gov.au/
\end{itemize}
